\documentclass{article}
\usepackage{graphicx, hyperref} % Required for inserting images
\usepackage[dvipsnames]{xcolor}
\usepackage{amssymb}
\usepackage{amsmath}
\usepackage{amsthm}
\usepackage{bbm}
\newtheorem{theorem}{Theorem}
\newtheorem{lemma}{Lemma}
\newtheorem{definition}{Definition}

\usepackage{tikz}
\usepackage{url}

\tikzset{%
  zeroarrow/.style = {-stealth,dashed},
  onearrow/.style = {-stealth,solid},
  c/.style = {circle,draw,solid,minimum width=2em,
        minimum height=2em},
  r/.style = {rectangle,draw,solid,minimum width=2em,
        minimum height=2em}
}

\newcommand{\midlinewidth}{1.0pt}
\newcommand{\incmidlinewidth}{1.7pt}
\newcommand{\middist}{20pt}
\newcommand{\halfdist}{13pt}
\definecolor{petroil2} {RGB} {36, 165, 175}
\definecolor{gold2} {RGB} {255, 130, 0}

\title{Causal Explanations in Chirho}
\author{}
\date{\today}

\newcommand{\poorva}[1]{{\color{blue}[Poorva: #1]}}

\begin{document}

\maketitle

% SAM's NOTES


\section{Intended audience}

General ML community with some interest in explainability and causal modeling, but not just people deep into causal inference, and definitely not philosophers.


\section{Tasks}


POTENTIAL:

    ---
    SUSPECTS  A  -> S 
    
    ----
    
    WITNESSES  W
    -----

ACTIVE


    ---
    ACTIVE SUSPECTS    C
    ACTIVE WITNESSES    T


bold face instead of vectors; S, W, C, T, bold face
remove arrows




2. Decide on terminology, make sure it's used consistently throughout 
Frontload notation
{\color{red}Dan: write notation section. Use boldface for sets and vectors. Keep a dictionary of terms and notation.}

3. Finish section on the continuous case (revise notation/terminology as necessary), making sure there is a last paragraph describing practical considerations 

3.5 Decide on whether we want to generalize/upgrade the definitions for the categorical case

4. Make the discussion about theorems a bit more exciting, and move proofs to an appendix. Make sure the proofs in the AC section are correct, and finish the last proof. Make sure the section looks relatively complete

5. Write up HouseIQ case study

6. Write up SHAP section (optional)

MILESTONE








\section{Claims}

NAME is a meta-modeling procedure that extends probabilistic SCM with latent "intervention" and "witness" auxiliary variables representing whether causal interventions are applied and what part of the context is kept fixed, with ... [list all features]

\begin{enumerate}
    \item \label{claim:unification} \textbf{Theoretical Unification}: NAME provides a principled causal probabilistic toolkit inspired by previously disparate notions in actual causality, causal explanation, etc., all as particular instantiations of probabilistic inference in an auxiliary variable model.
    
    \item \label{claim:algorithmic} \textbf{Algorithmic}: queries using NAME can be evaluated (or approximated) using existing off-the-shelf algorithmic and computational tools for probabilistic inference to start with, and using or implementing state-of-the-art algorithmic methods implemented in modern hardware at little additional cost - 
    \item \label{claim:extensions} \textbf{Extensions}: NAME follows a fairly natural path of going full Bayesian about elements already present in the existing notions in the vicinity
    
    \item It goes beyond what the existing tools do in:
        \begin{enumerate}
            \item \label{claim:ex-approximate} approximate search for explanation via approximate probabilistic inference,
            \item \label{claim:ex-continuous} allow continuous and categorical variables as inputs and outputs
            \item \label{claim:ex-probabilistic} replaces combinatorial and computationally hard components with probabilistic components, which can be feasibly/rigorously approximated
        \end{enumerate}
\end{enumerate}


\section{Evidence}

\begin{itemize}

    \item Conceptual discussion showing how this approach extends previous attempts, and justifying it as a reasonable thing to do. Almost a philosophical argument. Supports~\ref{claim:unification} and \ref{claim:extensions}. Include a look at Pearl's probability of causation from the book
    
    \item Theorems showing exact equivalence between PON and POS and probabilistic queries in our approach. Supports~\ref{claim:unification}. (no witnesses? )
    
    \item Theorems showing exact equivalence between other actual causality queries \textbf{not including minimality} and probabilistic queries in our approach. Supports~\ref{claim:unification}.
    
%    \item TBD Theorems about LIME and/or Shapley. Somewhat speculative. Supports~\ref{claim:unification}.
    
    \item Case studies of TBD. Supports~\ref{claim:unification}, \ref{claim:algorithmic}, and \ref{claim:extensions}.
    
    \item Implementation in ChiRho, that uses probabilistic inference from Pyro. Supports~\ref{claim:algorithmic}

    \item Comparison to less intuitive results provided by Shap or LIME
\end{itemize}


% END SAM's NOTES

% Eli: emphasize unity
% Eli: Beckers: https://link.springer.com/article/10.1007/s10992-021-09601-z

% Eli: psychological experiment papers
% - Lagnado
% - Gerstenberg
% - Icard
% - Kleiman-Weiner

% Eli: related quantity, algorithmic recourse, 
% responsibility/blame
%https://arxiv.org/abs/2010.10596
%Counterfactual Explanations and Algorithmic Recourses %for Machine Learning: A Review
% AI safety/agency
%https://arxiv.org/abs/2402.07221
%The Reasons that Agents Act: Intention and Instrumental Goals
%https://aaai.org/papers/11557-towards-formal-definitions-of-blameworthiness-intention-and-moral-responsibility/
%https://www.ijcai.org/proceedings/2023/41
%ijcai.orgijcai.org
%Quantifying Harm | IJCAI
%Electronic proceedings of IJCAI 2023


\section{Open issues}

\begin{itemize}
\item nicer ML-oriented examples
\item Benchmark?? Database queries?
\item read more psychogical papers to figure an example more connected with what's already published and goes beyong the simple philosophical examples

\end{itemize}


\section{Introduction}

\begin{itemize}

\item Explanation and attribution, short handwavy intro

\item  Bayesian causal probabilistic programs Witty2023BayesianStructuralCausalInference

mention \url{https://github.com/tum-i4/causal-canvas}, ibrahim2020canvas


\end{itemize}


\section{Motivations: actual causality and related notions}


\begin{itemize}

\item Start with SCM (take from the math tex) and the forest fire example

\item but-for clause and its failure

% Some quantities of interest for the attribution problem  have already been suggested.  There are two key concepts on the market, both of them focusing on what would have happened if a certain event had been absent. One of them is that of It turns out that  this notion is not fine-grained enough to correctly answer attribution queries in the context of complex causal models that involve over-determination and causal undercutting, so another notion, that  of \emph{actual causality}, has been proposed.\cite{halpern2016actual} 

\item actual causality

\begin{itemize}

\item Paying attention to contexts temporal2010kleinberg

\item Psychological validity blame2010gerstenbergLagnado

\item Cost minimization perspective cost2017zhou

\item formal properties: intransitivty and assymetry transitivity2017beckers.

\item computational considerations for exact computations: known complexity results complexityCausality2002Eiter, Aleksandrowicz2014computational,  checking2020ibrahim, tractable cases causesExplanations2002Eiter and tree pruning Hopkins2002strategies


\end{itemize}


\item Probability of attribution necessity/sufficiency/causation/attribution, robins1989probability, pearl2022probabilities, rosenbaum2002attributing,yamamoto2012understanding which is in some sense a causal cousin of quantities known in health sciences and excessive risk ratio and observational risk ratio. corradi2020causes



\item Individual causal effects individual2013kenneth individualized2021Hoogland, post2021individual, individual2013kenneth




\end{itemize}



\section{Causal explanation: informal explanation}

\begin{itemize}

\item How we depart/generalize 


This needs more thinking/items/ordering



\begin{itemize}

\item Going probabilistic proposed2015fenton-gryll



\item Probabilistic search through potential cause sets and potential alternative values

%As variables can interact and we want our methods be sensitive to potential systematic interventions, the tool needs to search not only through upstream variables but also sets thereof. Such variables potentially are not binary, so we should also search through their alternative values as well. As this is a large space, the only viable approach is to employ an approximate stochastic search.



\item Allowing continuous outcome and input variables

%Often a researcher is interested also in continuous outcomes (say, high quality years of life) or interventions (such as dosage or exposure level). Methods to search for causes in such a setting would be valuable.





\item Joining necessity and sufficiency   sufficiency2021beckers pearl2022probabilities

% Often the researchers care about more than just necessary causes: someone being alive is a necessary cause for their illness, but it is never an interesting cause to report. For this reason we should in parallel test for probabilistic sufficiency of a given cause (set), employing a notion akin to that of probability of necessity and sufficiency.

\item Controllable level of individualization: 
%The question of what would happen if a certain variable had a different value is ambiguous in many ways, some of which (as it transpired from the discussions surrounding actual causality) relate to which of other variables do we prefer to hold (potentially stochastically) fixed---the researcher should have a gradable control over such factors.

\item Going soft on equality and non-equality tavares19predicateExchange

\end{itemize}

\item  Informal gloss on what's going on in the next section

\end{itemize}



\section{Probabilistic explanation}


This is where the core content comes in.


\section{Examples}

\subsection{Over-determination: forest fire}

As in notebook

\subsection{Over-determination and under-cutting: stone throwing}

As in the SIR notebook

\subsection{Search for causes: complex forest fire}



\section{Connection to actual causality}

One-way theorem (see Poorva's suggestion in main.tex)


\section{Connection to explainable AI (potentially another paper)}

\begin{itemize}

\item In a degenerate setting: Shapley 

\item Others? Potentially in a different paper.

\end{itemize}

\section{Downstream applications}

Brief mention, no extended discussion

\begin{itemize}
\item  causal perspective on representation learning, formalizing non-spuriousness and efficiency in supervised representation learning, Wang2021desiderataCausal


\item Causes take the form of database tuples, and
can be ranked according to their causal responsibility, a numerical
measure of their relevance as a cause for the query answer. Salimi2016databases Salimi2014causesToQueries databases2010meliou atribution2023Bertossi 

\item Protocols involving multiple agents and their interactions, blame assignments that captures agent interactions and agents’ choices to deviate. interaction2015Sharma.pdf Halpern2019blameworthiness.pdf, Ward2024reasonsAgentsAct.pdf

\item team plans as causal models, responsibility in such teams2020alechina.pdf

\item software dependence, automatic checks for form  of
dependence that were previously too abstract or high-level
to be candidates for tools. Koppel2020demystifyingDependence.pdf
\end{itemize}



\bibliographystyle{plain}
\bibliography{references}


\end{document}
